% ============================================================
%  Глава 4 — Реализация регулятора в роботе
% ============================================================
\chapter{Реализация в коде робота}
\label{ch:implementation}

В этой главе описывается \textbf{полный конвейер}:
\textsf{MATLAB (offline)} $\to$ \textsf{YAML/JSON} $\to$ \textsf{Python (online)}
$\to$ \textsf{motor\_node} реального робота. Все коэффициенты выносятся
в \filepath{config.yaml}, что позволяет менять поведение робота
без перекомпиляции.

\section{Архитектура}
\label{sec:architecture}

\begin{center}
\begin{tikzpicture}[
  box/.style={draw, rounded corners, minimum width=3cm, minimum height=0.9cm,
              font=\small, align=center},
  arrow/.style={-{Stealth[length=2.5mm]}, thick},
  node distance=0.6cm and 1.5cm
]
\node[box, fill=blue!8]   (matlab)  {MATLAB\\\filepath{matlab/main.m}};
\node[box, fill=green!8, right=of matlab] (yaml) {\filepath{config.yaml}\\\small \texttt{control:} раздел};
\node[box, fill=orange!8, right=of yaml] (py)   {Python\\\filepath{pi\_nodes/control/}};
\node[box, fill=red!8, below=of py] (motor) {\filepath{motor\_node.py}\\цикл 20 Гц};

\draw[arrow] (matlab) -- node[above, font=\scriptsize] {export} (yaml);
\draw[arrow] (yaml)   -- node[above, font=\scriptsize] {load} (py);
\draw[arrow] (py)     -- node[right, font=\scriptsize] {cmd\_vel} (motor);
\end{tikzpicture}
\end{center}

\textbf{Шаги работы:}
\begin{enumerate}
  \item Инженер запускает \texttt{main.m} в MATLAB --- получает $\mat{A}, \mat{B}, \mat{K}, \mat{L}, \mat{P}_f$;
  \item Скрипт автоматически экспортирует матрицы в \filepath{config.yaml}
        раздел \texttt{control};
  \item Робот при старте читает \texttt{control} и инициализирует регулятор;
  \item На каждом цикле \texttt{motor\_node} вызывает $\vect{u}_k = \texttt{controller.step}(\vect{x}_k, \vect{x}_k^{\text{ref}})$;
  \item Если \texttt{control.mode} = \texttt{off}, действует старая логика
        прямой трансляции \texttt{cmd\_vel}.
\end{enumerate}

\section{Конфигурация}
\label{sec:config}

Все коэффициенты вынесены в \filepath{config.yaml}, раздел \texttt{control}:

\begin{filebox}[title={\filepath{config.yaml} --- \texttt{control:}}]
\begin{lstlisting}[basicstyle=\footnotesize\ttfamily]
control:
  mode: "off"           # off | lqr | mpc | modal — режим регулятора

  plant:                # параметры объекта (физика)
    v0:        0.20     # м/с — опорная скорость линеаризации
    tau_v:     0.15     # с — постоянная мотора (linear)
    tau_w:     0.10     # с — постоянная мотора (angular)
    Ts:        0.05     # с — период дискретизации (= 1 / 20 Hz)

  weights:              # критерий качества
    Q_diag:  [10, 10, 5, 1, 1]   # x,y,theta,v,omega
    R_diag:  [1, 1]              # u_v, u_omega

  limits:               # физические ограничения управления
    u_min:   [-0.30, -2.0]
    u_max:   [+0.30, +2.0]

  mpc:                  # параметры MPC
    horizon_N:  10      # горизонт прогноза (шагов)
    solver:     "clip"  # clip | qp (qp требует scipy/cvxpy)

  modal:                # параметры модального регулятора
    poles_continuous: [-2, -3, -4, -8, -10]   # желаемые
    observer_speedup: 5.0   # во сколько раз быстрее регулятора

  # Матрицы рассчитаны MATLAB-скриптом, обновляются автоматически.
  # Для ручной правки: установите override: true и поменяйте значения.
  matrices:
    override: false
    A: null   # 5x5 — заполнится из MATLAB
    B: null   # 5x2
    K: null   # 2x5 — главный регулятор
    L: null   # 5x5 — наблюдатель (если включён)
    Pf: null  # 5x5 — терминальный штраф MPC
\end{lstlisting}
\end{filebox}

\section{Python-модуль управления}
\label{sec:python_impl}

В пакете \filepath{pi\_nodes/control/} добавлены три новых модуля:

\begin{center}
\begin{tabular}{ll}
\toprule
\textbf{Файл} & \textbf{Назначение}\\
\midrule
\filepath{state\_space\_model.py} & Дискретная модель $\vect{x}_{k+1}=\mat{A}_d\vect{x}_k + \mat{B}_d\vect{u}_k$\\
\filepath{lqr\_controller.py}     & ЛКР: $\vect{u}_k = -\mat{K}\,\vect{x}_k$\\
\filepath{mpc\_controller.py}     & MPC: явное решение + проекция на бокс\\
\bottomrule
\end{tabular}
\end{center}

\subsection{Загрузка матриц}

\begin{filebox}[title={\filepath{pi\_nodes/control/state\_space\_model.py}}]
\begin{lstlisting}[style=pythonstyle]
import numpy as np
from config_loader import cfg

class StateSpaceModel:
    """Discrete state-space model x[k+1] = Ad*x[k] + Bd*u[k]."""
    def __init__(self, Ad=None, Bd=None):
        if Ad is None:
            Ad = np.array(cfg('control.matrices.A'), dtype=float)
        if Bd is None:
            Bd = np.array(cfg('control.matrices.B'), dtype=float)
        self.Ad, self.Bd = Ad, Bd
        self.n, self.r = self.Bd.shape

    def step(self, x, u):
        return self.Ad @ x + self.Bd @ u
\end{lstlisting}
\end{filebox}

\subsection{Регулятор LQR}

\begin{filebox}[title={\filepath{pi\_nodes/control/lqr\_controller.py}}]
\begin{lstlisting}[style=pythonstyle]
import numpy as np
from config_loader import cfg

class LQRController:
    def __init__(self, K=None, u_min=None, u_max=None):
        self.K     = np.array(K     if K     is not None
                              else cfg('control.matrices.K'))
        self.u_min = np.array(u_min if u_min is not None
                              else cfg('control.limits.u_min'))
        self.u_max = np.array(u_max if u_max is not None
                              else cfg('control.limits.u_max'))

    def step(self, x, x_ref=None):
        dx = x - x_ref if x_ref is not None else x
        u  = -self.K @ dx
        return np.clip(u, self.u_min, self.u_max)
\end{lstlisting}
\end{filebox}

\subsection{MPC: явное решение + проекция}

\begin{filebox}[title={\filepath{pi\_nodes/control/mpc\_controller.py} (упрощённо)}]
\begin{lstlisting}[style=pythonstyle]
class MPCController:
    """MPC c явным решением + проекцией на бокс ограничений."""

    def __init__(self, Ad, Bd, Q, R, Pf, N, u_min, u_max):
        self.N = N
        self._build_qp_matrices(Ad, Bd, Q, R, Pf)
        self.u_min, self.u_max = u_min, u_max
        self.r = Bd.shape[1]

    def _build_qp_matrices(self, Ad, Bd, Q, R, Pf):
        n, r = Bd.shape
        N = self.N
        # Phi (n*N x n) и Gamma (n*N x r*N) -- развёртка динамики
        Phi   = np.vstack([np.linalg.matrix_power(Ad, i+1) for i in range(N)])
        Gamma = np.zeros((n*N, r*N))
        for i in range(N):
            for j in range(i+1):
                block = np.linalg.matrix_power(Ad, i-j) @ Bd
                Gamma[i*n:(i+1)*n, j*r:(j+1)*r] = block

        Q_bar = np.kron(np.eye(N-1), Q)
        Q_bar = block_diag_pad(Q_bar, Pf)            # терминальный штраф
        R_bar = np.kron(np.eye(N), R)

        H_full = 2*(Gamma.T @ Q_bar @ Gamma + R_bar)
        f_mat  = 2*Gamma.T @ Q_bar @ Phi

        # Только первое управление: K_first = (H^-1 f)[:r, :]
        K_first = -np.linalg.solve(H_full, f_mat)[:r, :]
        self.K_first = K_first

    def step(self, x, x_ref=None):
        dx = x - x_ref if x_ref is not None else x
        u  = self.K_first @ dx
        return np.clip(u, self.u_min, self.u_max)
\end{lstlisting}
\end{filebox}

\section{Тестирование}
\label{sec:testing}

Для каждого модуля написан unit-тест в \filepath{tests/}.

\subsection{test\_state\_space\_model.py}

Проверяет:
\begin{itemize}
  \item $\vect{x}_{k+1} = \mat{A}_d\vect{x}_k + \mat{B}_d\vect{u}_k$ ---
        корректность дискретного шага;
  \item Сходимость к точке равновесия при $\vect{u} = \vect{0}$
        (для устойчивого $\mat{A}_d$).
\end{itemize}

\subsection{test\_lqr\_controller.py}

Проверяет:
\begin{itemize}
  \item Замкнутая система $\mat{A}_d - \mat{B}_d\mat{K}$ устойчива
        ($|\lambda_i| < 1$);
  \item Управление лежит в пределах $[\vect{u}_{\min}, \vect{u}_{\max}]$;
  \item При $\vect{x} = \vect{0}$ управление равно нулю.
\end{itemize}

\subsection{test\_mpc\_controller.py}

Проверяет:
\begin{itemize}
  \item Без ограничений MPC совпадает с ЛКР для горизонта $N \geq 20$;
  \item При активных ограничениях управление лежит на границе бокса;
  \item Замкнутая симуляция сходится к нулю.
\end{itemize}

\section{Примеры использования}
\label{sec:examples}

\subsection{Запуск симуляции с MPC}

\begin{filebox}[title={Изменение режима регулятора}]
\begin{lstlisting}[basicstyle=\small\ttfamily]
# 1. В config.yaml:
control:
  mode: mpc

# 2. Запустить симулятор:
./samurai.sh sim

# 3. Открыть dashboard, отправить goal_pose
\end{lstlisting}
\end{filebox}

\subsection{Расчёт нового регулятора в MATLAB}

\begin{filebox}[title={\filepath{matlab/main.m} --- сценарий}]
\begin{lstlisting}[style=matlabstyle]
% Загрузить параметры объекта
params = samurai_params();
% Линеаризовать
[A, B, C, D] = linearize_samurai(params);
% Дискретизировать
[Ad, Bd]     = discretize_samurai(A, B, params.Ts);
% Синтез ЛКР
[K_lqr, P]   = design_lqr(Ad, Bd, params.Q, params.R);
% Синтез MPC
mpc          = design_mpc(Ad, Bd, params.Q, params.R, P, params.N);
% Экспорт в config.yaml
export_to_yaml(Ad, Bd, K_lqr, P, mpc, '../config.yaml');
% Симуляция
simulate_closed_loop(Ad, Bd, K_lqr);
\end{lstlisting}
\end{filebox}

\subsection{«Ошибочный» режим для демонстрации}

Чтобы показать студентам \textbf{неустойчивую систему}, достаточно
поставить в \texttt{config.yaml}:
\begin{filebox}[title={Демонстрация неустойчивости}]
\begin{lstlisting}[basicstyle=\small\ttfamily]
control:
  mode: modal
  modal:
    poles_continuous: [+1.0, +0.5, -1, -1, -1]   # 2 положительных!
\end{lstlisting}
\end{filebox}

При запуске симуляции робот будет «убегать» по экспоненте,
так как замкнутая система имеет полюса в правой полуплоскости.

\section{Заключение}
\label{sec:conclusion}

\textbf{Резюме созданной системы:}

\begin{enumerate}
  \item Построена 5-мерная нелинейная ОДУ-модель робота
        (глава~\ref{ch:model});
  \item Модель линеаризована вокруг режима прямолинейного движения
        и дискретизирована с $T_s = 0{,}05$ с;
  \item Синтезированы 3 регулятора: модальный (pole placement),
        ЛКР (Risk-метод), MPC с горизонтом $N = 10$ и проекцией;
  \item Полноразмерный наблюдатель восстанавливает все 5 координат
        состояния по доступным измерениям;
  \item Расчёт выполняется offline в MATLAB
        (\filepath{matlab/}, 6 файлов), результат экспортируется в
        \filepath{config.yaml};
  \item Python-реализация (\filepath{pi\_nodes/control/}, 3 файла)
        читает конфиг при старте и применяет регулятор в цикле
        \texttt{motor\_node} (20 Гц);
  \item Покрыто unit-тестами (\filepath{tests/test\_*\_controller.py}, 3 файла);
  \item Все параметры $\mat{Q}, \mat{R}, N, \vect{u}_{\min,\max}$ ---
        редактируемы из конфига без перекомпиляции, что позволяет как
        улучшать поведение робота, так и намеренно вводить ошибки
        для отладки и обучения.
\end{enumerate}

\textbf{Ссылки на файлы проекта:}

\begin{center}
\begin{tabular}{ll}
\toprule
\textbf{Файл} & \textbf{Содержимое}\\
\midrule
\filepath{matlab/main.m}                                & оркестратор расчёта\\
\filepath{matlab/samurai\_params.m}                     & параметры объекта\\
\filepath{matlab/linearize\_samurai.m}                  & матрицы $\mat{A}, \mat{B}$\\
\filepath{matlab/discretize\_samurai.m}                 & ZOH-дискретизация\\
\filepath{matlab/design\_lqr.m}                         & решение DARE\\
\filepath{matlab/design\_mpc.m}                         & MPC-матрицы\\
\filepath{matlab/simulate\_closed\_loop.m}              & симуляция\\
\filepath{matlab/export\_to\_yaml.m}                    & запись конфига\\
\filepath{pi\_nodes/control/state\_space\_model.py}     & дискретная модель\\
\filepath{pi\_nodes/control/lqr\_controller.py}         & ЛКР-регулятор\\
\filepath{pi\_nodes/control/mpc\_controller.py}         & MPC-регулятор\\
\filepath{config.yaml}                                  & раздел \texttt{control:}\\
\filepath{tests/test\_state\_space\_model.py}           & тест модели\\
\filepath{tests/test\_lqr\_controller.py}               & тест ЛКР\\
\filepath{tests/test\_mpc\_controller.py}               & тест MPC\\
\bottomrule
\end{tabular}
\end{center}
